% 有向环与有向无环图：长度为2的有向环；三角形的循环定向；三角形的无环定向1->2,1->3,2->3。
\begin{tikzpicture}[>=Stealth,
  v/.style={circle,draw=FigureStroke,fill=FigureBlueFill,line width=.8pt,minimum size=6.5mm,inner sep=0pt},
  a/.style={draw=FigureStroke,line width=.8pt,->,shorten >=1pt},
  ar/.style={draw=FigureRed,line width=1.2pt,->,shorten >=1pt}]
\begin{scope}
  \node[v] (u) at (0,0) {$i$}; \node[v] (w) at (1.5,0) {$j$};
  \draw[ar] (u) to[bend left=30] (w); \draw[ar] (w) to[bend left=30] (u);
  \node[align=center] at (.75,-1.35) {(a) $i\to j\to i$\\长度$2$的有向环};
\end{scope}
\begin{scope}[xshift=3.6cm]
  \node[v] (1) at (0,.75) {$1$}; \node[v] (2) at (-.85,-.6) {$2$}; \node[v] (3) at (.85,-.6) {$3$};
  \draw[ar] (1)--(2); \draw[ar] (2)--(3); \draw[ar] (3)--(1);
  \node[align=center] at (0,-1.35) {(b) $1\to2\to3\to1$\\含有向环};
\end{scope}
\begin{scope}[xshift=6.9cm]
  \node[v] (1) at (0,.75) {$1$}; \node[v] (2) at (-.85,-.6) {$2$}; \node[v] (3) at (.85,-.6) {$3$};
  \draw[a] (1)--(2); \draw[a] (2)--(3); \draw[a] (1)--(3);
  \node[align=center] at (0,-1.35) {(c) 无环定向\\底层仍是三角形};
\end{scope}
\end{tikzpicture}
